\section{Chứng minh phi tương tác và phép biến đổi Fiat–Shamir}
(Non-Interactive Proofs and Fiat-Shamir)

\subsection{(1) [20 điểm] Phép biến đổi Fiat–Shamir (Fiat–Shamir Transform)}

\subsubsection{(a) [10 điểm] Từ tương tác sang phi tương tác (From Interactive to Non-Interactive)}

Bạn đang chuyển đổi giao thức nhận dạng Schnorr thành một sơ đồ chữ ký số bằng phép biến đổi Fiat–Shamir.

\begin{center}
    \textbf{Trả lời:}
\end{center}

\textbf{i. Trình bày toàn bộ thuật toán ký cho thông điệp}

``Transfer \$1000 to Bob''

Bao gồm cách tính thách thức $c = H(A, Y, m)$.

\textbf{Bài làm:}

Để chuyển đổi giao thức Schnorr tương tác thành chữ ký số phi tương tác (Non-Interactive Digital Signature), ta áp dụng phép biến đổi Fiat-Shamir. Ý tưởng cốt lõi là thay thế vai trò của Verifier (người gửi thách thức ngẫu nhiên) bằng một hàm băm mật mã (Hash Function).

**Thuật toán ký (Signing Algorithm):**
Giả sử Alice có cặp khóa $(a, A=g^a)$ và muốn ký lên thông điệp $m = \text{``Transfer \$1000 to Bob''}$.

\begin{enumerate}
    \item \textbf{Bước 1: Tạo cam kết (Commitment)}
    Alice chọn một số ngẫu nhiên $y$ (nonce).
    Tính cam kết $Y = g^y \pmod{n}$.
    
    \item \textbf{Bước 2: Tự tính thách thức (Compute Challenge)}
    Thay vì đợi Verifier gửi $c$, Alice tự tính $c$ bằng cách băm tất cả các dữ liệu công khai liên quan đến phiên làm việc này:
    \[
    c = H(A, Y, m) = H(A, Y, \text{``Transfer \$1000 to Bob''})
    \]
    Việc này đảm bảo $c$ phụ thuộc chặt chẽ vào thông điệp và cam kết, không thể thay đổi tùy tiện.
    
    \item \textbf{Bước 3: Tính phản hồi (Compute Response)}
    Alice tính phản hồi $r$ giống như trong giao thức tương tác:
    \[
    r = y + c \cdot a \pmod{n}
    \]
    
    \item \textbf{Bước 4: Xuất chữ ký (Output Signature)}
    Chữ ký cuối cùng cho thông điệp $m$ là cặp giá trị $\sigma = (c, r)$ (hoặc đôi khi là $(Y, r)$ tùy biến thể, nhưng $(c, r)$ phổ biến hơn vì kích thước nhỏ gọn).
\end{enumerate}

\textbf{ii. Giải thích vì sao việc đưa thông điệp $m$ vào trong hàm băm là cực kỳ quan trọng}

Nếu thay vào đó sử dụng $c = H(A, Y)$, thì loại tấn công nào sẽ trở nên khả thi?

\textbf{Bài làm:}

Việc đưa thông điệp $m$ vào đầu vào của hàm băm là **bắt buộc** để liên kết chữ ký với nội dung văn bản. Nếu ta loại bỏ $m$ và chỉ tính $c = H(A, Y)$, hệ thống sẽ dính lỗ hổng nghiêm trọng gọi là **Tấn công giả mạo sự tồn tại (Existential Forgery Attack)**.

**Kịch bản tấn công:**
\begin{itemize}
    \item Giả sử Alice ký một thông điệp vô hại $m_1 = \text{``Hello''}$. Cô ấy tạo ra chữ ký $(c, r)$ với $c = H(A, Y)$.
    \item Kẻ tấn công (Eve) bắt được chữ ký $(c, r)$ này.
    \item Eve muốn giả mạo chữ ký của Alice cho thông điệp $m_2 = \text{``Transfer \$1M to Eve''}$.
    \item Vì quy trình kiểm tra chữ ký chỉ là: $c \stackrel{?}{=} H(A, Y)$ và $g^r \stackrel{?}{=} Y \cdot A^c$.
    \item Eve nhận thấy rằng việc kiểm tra $c$ hoàn toàn không phụ thuộc vào $m$. Do đó, chữ ký $(c, r)$ hợp lệ cho $m_1$ cũng sẽ **hợp lệ cho bất kỳ thông điệp $m_2$ nào khác**.
    \item Eve chỉ cần copy-paste chữ ký của $m_1$ sang cho $m_2$, và Verifier sẽ chấp nhận nó.
\end{itemize}
Do đó, $m$ phải nằm trong hàm băm để mỗi chữ ký là duy nhất cho một nội dung cụ thể.

\vspace{0.5cm}

\subsubsection{(b) [10 điểm] Bảo mật trong mô hình Random Oracle}

Phép biến đổi Fiat–Shamir giả định rằng hàm băm hoạt động như một Random Oracle.

\begin{center}
    \textbf{Trả lời:}
\end{center}

\textbf{i. Giải thích những thuộc tính của Random Oracle mà một hàm băm thực tế (như SHA-256) có thể không đạt được hoàn hảo}

\textbf{Bài làm:}

**Random Oracle (RO)** là một mô hình lý tưởng hóa trong toán học, giống như một "hộp đen" hoàn hảo:
\begin{itemize}
    \item Với mỗi đầu vào mới, nó trả về một đầu ra ngẫu nhiên hoàn toàn (uniform random) từ không gian đích.
    \item Nếu gặp lại đầu vào cũ, nó trả về đúng kết quả đã tạo ra trước đó.
    \item Không ai có thể đoán trước hay tính toán đầu ra nếu không thực sự gọi hàm.
\end{itemize}

**Hàm băm thực tế (như SHA-256)** khác biệt ở chỗ:
\begin{itemize}
    \item **Tính xác định (Deterministic) và có cấu trúc:** SHA-256 là một thuật toán cụ thể, không phải ngẫu nhiên. Nếu bạn phân tích kỹ thuật toán, bạn có thể tìm thấy các mối quan hệ toán học giữa đầu vào và đầu ra (dù rất khó).
    \item **Length Extension Attack:** Với các hàm băm cấu trúc Merkle-Damg\aa rd (như SHA-256), nếu biết $H(x)$, ta có thể tính được $H(x || \text{padding} || y)$ mà không cần biết $x$. Random Oracle không cho phép điều này (đầu ra của chuỗi mới phải hoàn toàn ngẫu nhiên và không liên quan đến chuỗi cũ).
    \item **Khả năng bị tấn công:** Trong thực tế, đã có những trường hợp (như tấn công vào Fiat-Shamir khi dùng hàm băm yếu hoặc thiết kế sai) mà kẻ tấn công lợi dụng cấu trúc của hàm băm để tạo ra bằng chứng giả, điều không thể xảy ra với Random Oracle.
\end{itemize}

\textbf{ii. Mô tả ở mức cao Bổ đề Forking (Forking Lemma)}

Giải thích cách bổ đề này cho phép rút trích nhân chứng (witness) từ một kẻ tấn công có khả năng giả mạo chứng minh Fiat–Shamir.

\textbf{Bài làm:}

**Bổ đề Forking (Forking Lemma)** là công cụ lý thuyết quan trọng để chứng minh độ an toàn của chữ ký Fiat-Shamir. Nó hoạt động dựa trên kỹ thuật "tua ngược thời gian" (rewinding) trong máy ảo.

**Cơ chế:**
\begin{itemize}
    \item Giả sử có một kẻ tấn công (Adversary) có thể tạo ra chữ ký giả mạo $(Y, c, r)$ với xác suất không nhỏ.
    \item Ta chạy kẻ tấn công lần đầu tiên. Khi nó hỏi Random Oracle để lấy giá trị băm cho $(A, Y, m)$, ta trả về thách thức $c$. Kẻ tấn công xuất ra chữ ký hợp lệ $(c, r)$.
    \item Sau đó, ta **tua ngược** (rewind) kẻ tấn công lại đúng thời điểm nó vừa gửi truy vấn băm đó. Mọi trạng thái bộ nhớ của nó được giữ nguyên như cũ.
    \item Lần này, ta trả về một thách thức **khác** $c'$ (ngẫu nhiên mới). Do trạng thái nội tại giống hệt lần trước, kẻ tấn công (nếu đủ giỏi) vẫn sẽ tạo ra được chữ ký hợp lệ thứ hai $(c', r')$ với cùng cam kết $Y$.
    \item Kết quả: Ta thu được hai bộ $(Y, c, r)$ và $(Y, c', r')$. Đây chính là điều kiện để áp dụng kỹ thuật trích xuất (như đã làm ở bài 1) để tìm ra khóa bí mật $a$.
\end{itemize}
Tóm lại, Bổ đề Forking cho thấy: Nếu ai đó giả mạo được chữ ký, họ thực sự phải biết khóa bí mật.

\vspace{0.5cm}

\subsection{(2) [15 điểm] Mạch và Zero-Knowledge tổng quát (Circuits and General-Purpose ZK)}

Bạn muốn chứng minh rằng mình biết một giá trị $x$ sao cho
$\text{SHA-256}(x) = y$
với một $y$ đã cho.

\begin{center}
    \textbf{Trả lời:}
\end{center}

\textbf{(a) [7.5 điểm] Giải thích vì sao việc tạo một giao thức Sigma riêng cho phát biểu này là rất khó khăn}

Làm thế nào việc biểu diễn SHA-256 dưới dạng mạch (circuit) lại khiến chứng minh này trở nên khả thi?

\textbf{Bài làm:}

\begin{itemize}
    \item **Khó khăn của Sigma Protocol:** Giao thức Sigma (như Schnorr) dựa trên các tính chất đại số đẹp đẽ của lý thuyết số (như $g^a \cdot g^b = g^{a+b}$). Trong khi đó, SHA-256 là một "mớ hỗn độn" các phép toán bitwise phi tuyến tính (XOR, AND, Rotate, Shift) được thiết kế để phá vỡ mọi cấu trúc đại số. Việc cố gắng xây dựng một phương trình đại số $Y = g^x$ tương đương cho hàng ngàn vòng lặp và phép toán bit của SHA-256 là cực kỳ phức tạp và gần như bất khả thi để làm thủ công.
    
    \item **Giải pháp Mạch (Circuits):** Thay vì tìm công thức toán học, ta "vẽ" lại toàn bộ quy trình tính toán của SHA-256 dưới dạng một mạch điện tử khổng lồ (Arithmetic Circuit hoặc Boolean Circuit). Mạch này gồm hàng ngàn cổng logic (cổng cộng, cổng nhân, cổng XOR).
    \item Các hệ thống ZK tổng quát (General-purpose ZK) như zk-SNARKs (sử dụng R1CS hoặc QAP) cho phép ta chứng minh rằng: "Tôi biết các tín hiệu đi vào dây dẫn sao cho tất cả các cổng trong mạch đều đúng". Máy tính có thể tự động chuyển đổi mạch SHA-256 thành hệ thống ràng buộc đa thức này mà không cần con người phải thiết kế giao thức riêng.
\end{itemize}

\textbf{(b) [7.5 điểm] Ứng dụng của bạn cần chứng minh mệnh đề}

``Tôi biết $x$ sao cho $x^3 + x + 5 = 35$ trong $\mathbb{Z}_{101}$''

Bạn sẽ chọn giao thức tùy chỉnh hay dạng mạch tổng quát?
Hãy giải thích lựa chọn của bạn.

\textbf{Bài làm:}

Tôi sẽ chọn **Giao thức tùy chỉnh (Custom Protocol)** hoặc một biến thể Sigma Protocol đơn giản cho đa thức.

**Lý do:**
\begin{itemize}
    \item **Hiệu suất:** Bài toán $x^3 + x - 30 = 0$ là một quan hệ đại số cực kỳ đơn giản và nhỏ gọn.
    \item **Overhead của Mạch tổng quát:** Nếu dùng zk-SNARKs (mạch tổng quát), ta phải chịu chi phí khởi tạo lớn (Trusted Setup, sinh Key), kích thước thư viện lớn, và thời gian tạo bằng chứng (Prover time) lâu hơn đáng kể do phải chuyển đổi đa thức đơn giản này thành hàng ngàn ràng buộc R1CS không cần thiết. Nó giống như việc dùng dao mổ trâu để giết gà.
    \item **Tính đơn giản:** Một giao thức Sigma chuyên biệt cho việc chứng minh kiến thức nghiệm của đa thức (dựa trên cam kết Pedersen chẳng hạn) sẽ rất nhanh, nhẹ, và dễ triển khai hơn nhiều cho trường hợp cụ thể này.
\end{itemize}
